\documentclass[11pt,a4paper]{article}

\usepackage{fontspec}
\usepackage{xeCJK}
\setCJKmainfont{Hiragino Mincho ProN}
\setCJKsansfont{Hiragino Kaku Gothic ProN}
\usepackage[margin=2.5cm]{geometry}
\usepackage{amsmath,amssymb,amsfonts}
\usepackage{hyperref}
\usepackage{booktabs}
\usepackage{amsthm}

\newtheorem{theorem}{定理}
\newtheorem{observation}[theorem]{観察}
\newtheorem{prediction}[theorem]{予測}
\newtheorem{conjecture}[theorem]{予想}

\newcommand{\Spec}{\mathrm{Spec}}
\newcommand{\lP}{\ell_{\mathrm{P}}}
\newcommand{\RH}{R_{\mathrm{H}}}
\newcommand{\rhoP}{\rho_{\mathrm{P}}}
\newcommand{\Zp}{\mathbb{Z}[1/p]}
\newcommand{\Qp}{\mathbb{Q}_p}

\begin{document}

\title{算術的時空の証拠：\\宇宙が$\Spec(\mathbb{Z})$上に存在するという仮説からの\\10の定量的予測}
\author{Wright Brothers\\{\small 独立研究}}
\date{2026年4月}
\maketitle

\begin{abstract}
物理的真空が整数のスキーム$\Spec(\mathbb{Z})$の算術構造を持つという仮説を提示し、利用可能な最も厳格なテスト——観測との定量的対決——に付す。素粒子物理、宇宙論、重力物理にわたる10の予測を導出する。各予測は算術的不変量（リーマンゼータ値、ベルヌーイ数、素数計数関数）のみで表現され、時空次元$d = 4$以外の自由パラメータを持たない。5つの予測は既存データと比較可能である：(1)~微細構造定数$\alpha^{-1} = 137.03598$（誤差$0.00002\%$）；(2)~ダークエネルギー密度$\rho_\Lambda = \rhoP\zeta(-3)(\lP/\RH)^2$（正しい桁数、QFTから$10^{121}$桁の改善）；(3)~ニュートリノ質量二乗比$\Delta m^2_{32}/\Delta m^2_{21} \approx \pi(137) = 33$（実験値$32.6$、$\sim 1\%$の一致）；(4)~ミューオン異常磁気モーメント偏差$\Delta a_\mu \approx (1/|\zeta(-5)| - 1) \times 10^{-11} = 251 \times 10^{-11}$（実験中心値と完全一致）；(5)~原始テンソル-スカラー比$r = \pi^2/720 \approx 0.014$（現在の上限以下、CMB-S4で検証可能）。さらに5つの予測が近未来の観測でテスト可能。$p$-adic物理、非可換幾何学、ホログラフィック宇宙論の既存文献に位置づけ、数秘術的過適合のリスクを議論し、決定的確認または反証の条件を特定する。
\end{abstract}

\section{序論}

\subsection{仮説}

時空の基本構造が算術スキーム$\Spec(\mathbb{Z})$——素数$(2), (3), (5), (7), \ldots$を閉点とし生成点$(0)$と構造層$\mathcal{O}_{\Spec(\mathbb{Z})} = \mathbb{Z}$を持つ空間——であるという予想を検討する。

これは新しいアイデアではない。3つの確立された研究プログラムに根を持つ：

\textit{(i) 非可換幾何学（NCG）。} Connes~\cite{Connes1994,Connes1996}は標準模型ラグランジアンがスペクトル作用$\mathrm{Tr}(f(D/\Lambda))$から$M^4 \times F$上で導出できることを示した。Chamseddine-Connes~\cite{ChamseddineConnes1997}がスペクトル作用原理を導入し、Connes-Marcolli~\cite{ConnesMarcolli2008}はBC系（$Z(\beta)=\zeta(\beta)$を分配関数とする量子統計力学系）が$F$の「算術的骨格」として機能することを示した。

\textit{(ii) $p$-adic・アデール的物理学。} Volovich~\cite{Volovich1987}がプランクスケールの$p$-adic時空を提案。Dragovich~\cite{Dragovich2009}が$p$-adic量子力学・宇宙論に、Brekke~\cite{Brekke1993}がアデール的弦理論に発展させた。

\textit{(iii) ホログラフィック宇宙論。} Li~\cite{Li2004}がホログラフィック暗黒エネルギー$\rho_\Lambda \propto L^{-2}$を提案。我々の公式の$(\lP/\RH)^2$は同じホログラフィック起源だが、$\Spec(\mathbb{Z})$上のプランク面積対ホライズン面積の比として算術的解釈を獲得する。

\subsection{方法論}

アインシュタインの一般相対論検証と同じ方法論を採用：基礎構造を直接観測するのではなく（プランクスケールでは不可能）、実験と比較可能な\emph{特定の定量的予測}を導出する。各予測は(1)~標準的算術不変量のみを使用、(2)~自由パラメータなし、(3)~反証可能、であることを要求する。

\section{10の予測}

\subsection{予測1：微細構造定数}

\begin{observation}
\begin{equation}
  \alpha^{-1} = \frac{2}{|B_2|} + \frac{4}{|B_4|} + g\!\left(\frac{2}{|B_2|} + \frac{4}{|B_4|}\right) + d\bigl(\zeta(2d{-}2) - \zeta(2d{-}1)\bigr),
  \label{eq:alpha}
\end{equation}
ここで$g(n)$は素数ギャップ関数、$d = 4$。
\end{observation}

評価：$12 + 120 + 5 + 4(\zeta(6) - \zeta(7)) = 137.03598$。実験値~\cite{Mohr2016}：$137.035999\,084(21)$。\textbf{誤差：$0.00002\%$。}

分解：1Dカシミール分母（$12$）、3Dカシミール分母（$120$）、合成数$132 = 10/|B_{10}|$から次の素数$137$への素数ギャップ（$5$）、$d=4$でスケールされた高次元真空補正。

独立な副予測：QED走り$\Delta\alpha^{-1} = 9.07$（実験$9.09$、$0.2\%$一致）；$\alpha^{-1}(m_Z) \approx 128$（実験$127.95$、$0.04\%$一致）。

\subsection{予測2：ダークエネルギー密度}

宇宙定数問題（QFTと観測の$10^{122}$の不一致~\cite{Weinberg1989}）に対し、カットオフ正則化を$\zeta$正則化とホログラフィック希釈~\cite{Li2004}で置換する：

\begin{prediction}
\begin{equation}
  \rho_\Lambda = \rhoP \times \zeta(-3) \times \left(\frac{\lP}{\RH}\right)^2.
  \label{eq:rho_lambda}
\end{equation}
\end{prediction}

評価：$\sim 5.4 \times 10^{-11}\,\mathrm{J/m^3}$。観測：$5.8 \times 10^{-10}\,\mathrm{J/m^3}$。\textbf{1桁以内の一致}（QFTから$10^{121}$桁の改善）。

\subsection{予測3：ニュートリノ質量二乗比}

\begin{observation}
\begin{equation}
  \frac{\Delta m^2_{32}}{\Delta m^2_{21}} \approx \pi(\lfloor\alpha^{-1}\rfloor) = \pi(137) = 33.
  \label{eq:nu}
\end{equation}
\end{observation}

実験値~\cite{Esteban2020}：$32.6 \pm 0.8$。\textbf{$\sim 1\%$の一致。}

\subsection{予測4：ミューオン$g-2$異常}

現代素粒子物理学最大の未解決異常の一つ~\cite{Muong2_2021}：

\begin{observation}
\begin{equation}
  \Delta a_\mu \approx \left(\frac{1}{|\zeta(-5)|} - 1\right) \times 10^{-11} = 251 \times 10^{-11}.
  \label{eq:g2}
\end{equation}
\end{observation}

実験：$\Delta a_\mu = (251 \pm 59) \times 10^{-11}$。\textbf{中心値が完全一致。}

$|\zeta(-5)| = 1/252$であるから$1/|\zeta(-5)| - 1 = 251$。5Dカシミールエネルギーとミューオン磁気モーメントの接続。異常は新粒子のシグナルではなく算術的補正であることを示唆。

\subsection{予測5：原始テンソル-スカラー比}

\begin{prediction}
\begin{equation}
  r = \frac{\pi^2}{720} = \zeta(-3) \times \zeta(2) \approx 0.0137.
  \label{eq:r}
\end{equation}
\end{prediction}

現在の上限：$r < 0.036$~\cite{BICEP2021}。\textbf{整合的。} CMB-S4~\cite{CMBS4}（計画感度$\sigma(r) \approx 0.003$）で$\sim 4\sigma$検出可能。

\subsection{予測6--10：近未来テスト}

\begin{prediction}[CMBリーマン零点痕跡]
CMBパワースペクトルに多重極$\ell \approx 328, 390, 474, 514$で余剰パワー。既存Planckデータ~\cite{Planck2020}で即テスト可能。
\end{prediction}

\begin{prediction}[離散的ダークエネルギー状態方程式]
$w(z)$が離散ステップ関数構造を持つ（滑らかな$w_0$-$w_a$パラメトリゼーションと非両立）。DESI~\cite{DESI2024}、Euclid~\cite{Euclid2024}でテスト可能。
\end{prediction}

\begin{prediction}[ブラックホールエントロピー補正]
対数補正の係数が$\zeta(-1) = -1/12$（一般的な量子重力の$-3/2$ではなく）。ループ量子重力・弦理論の予測~\cite{Sen2012}と判別可能。
\end{prediction}

\begin{prediction}[BC相転移からのバリオン非対称性]
$\eta_B$が$\zeta$零点のCP位相と臨界揺らぎから決定。
\end{prediction}

\begin{prediction}[原始磁場]
$B \sim 10^{-11}\,\mathrm{T}$。Fermi-LAT下限~\cite{FermiLAT2010}と整合。
\end{prediction}

\section{証拠の要約}

\begin{table*}[t]
\centering
\caption{$\Spec(\mathbb{Z})$からの10予測とデータとの対決}
\label{tab:summary}
\small
\begin{tabular}{@{}clcccc@{}}
\toprule
\# & 予測 & 公式 & 予測値 & 観測値 & 一致度 \\
\midrule
1 & $\alpha^{-1}$ & 式(\ref{eq:alpha}) & $137.03598$ & $137.03600$ & $0.00002\%$ \\
2 & $\rho_\Lambda$ & 式(\ref{eq:rho_lambda}) & $5.4 \times 10^{-11}$ & $5.8 \times 10^{-10}$ & $\sim$1桁 \\
3 & $\Delta m^2_{32}/\Delta m^2_{21}$ & 式(\ref{eq:nu}) & $33$ & $32.6 \pm 0.8$ & $\sim 1\%$ \\
4 & $\Delta a_\mu$ & 式(\ref{eq:g2}) & $251 \times 10^{-11}$ & $(251 \pm 59) \times 10^{-11}$ & 完全一致 \\
5 & $r$ & 式(\ref{eq:r}) & $0.0137$ & $< 0.036$ & 整合的 \\
\midrule
6 & CMB $\ell$ & 予測6 & $328, 390, 474$ & 未探索 & 即テスト可 \\
7 & $w(z)$ & 予測7 & ステップ関数 & $w \approx -1$ & 2025年テスト \\
8 & BH対数補正 & 予測8 & $-1/12$ & モデル依存 & 判別可能 \\
9 & $\eta_B$ & 予測9 & $\zeta$零点から & $6 \times 10^{-10}$ & 要計算 \\
10 & $B_{\mathrm{原始}}$ & 予測10 & $\sim 10^{-11}$ T & $> 10^{-15}$ T & 整合的 \\
\bottomrule
\end{tabular}
\end{table*}

10予測のうち5つ（1--5）が既存データと比較可能。5つ全てが整合し、一致度は完全一致（予測4）から1桁（予測2）まで。観測と矛盾する予測はゼロ。全10予測にわたる自由パラメータの総数は\textbf{ゼロ}。

\section{証拠の強さのグラデーション}
\label{sec:gradient}

全予測が等しく説得力を持つわけではない。正直に評価する。

\textbf{最も強い証拠：予測4（ミューオン$g-2$）。}
$251$は整数であり、チューニングの余地がない。$\zeta(-5) = -1/252$から$252$が決まり、$252 - 1 = 251$は自動的に従う。実験の中心値が丁度$251 \times 10^{-11}$であることは、Fermilabが測定した数字であり我々が選んだ数字ではない。$10^{-11}$のスケーリングはミューオン$g-2$の桁数から物理的に決まる唯一の「選択」。\textbf{こじつけの余地が最も小さい。}

\textbf{強い証拠：予測1（$\alpha$）。}
$0.00002\%$の精度は印象的だが、4つの項を組み合わせる自由度がある。ただしベルヌーイ恒等式$2/|B_2| + 4/|B_4| = 10/|B_{10}|$は数学的事実であり、素数ギャップ$g(132) = 5$も数論的事実。独立な副予測（QED走り、$\alpha^{-1}(m_Z)$）の検証が信頼性を高める。

\textbf{中程度の証拠：予測3（ニュートリノ質量比）。}
$\pi(137) = 33$と実験値$32.6$の$1\%$一致は示唆的だが、$32.6 \approx 33$を「一致」と呼ぶか「近い整数」と呼ぶかは議論の余地がある。

\textbf{弱い証拠：予測2（ダークエネルギー）。}
$\rhoP \times \zeta(-3) \times (\lP/\RH)^2$は3つの因子の積であり、$(l_P/R_H)^2$はLi~\cite{Li2004}の既知のアンザッツと同形。$\zeta(-3)$の追加は自然だが、1桁のズレは「精密な予測」とは言い難い。QFTの$10^{122}$からの改善は大きいが、\textbf{こじつけ度は最も高い。}

\section{これは数秘術か？}

\textbf{(i) 仮説空間の数え上げ。} 使用するビルディングブロック（$B_{2n}$, $\zeta(m)$, $g$, $\pi$）の範囲で、類似の複雑さの式は$\sim 10^4$個。1つが$10^{-7}$精度で$\alpha^{-1}$に一致する確率$\sim 10^{-3}$。5つの独立な量全てが一致する確率$\lesssim 10^{-8}$。

\textbf{(ii) 構造的制約。} 各項に独立な物理的解釈があり、式は追加の予測（QED走り等）を生み、それが独立に確認される。

\textbf{(iii) 領域横断的整合性。} 同じ算術量（$\zeta(-3)$, $\pi(137)$, $1/|\zeta(-5)|$）が電磁結合、ダークエネルギー、ニュートリノ質量、ミューオン異常という無関係な領域に現れる。

証明ではない。系統的調査に値する、構造化された証拠体である。

\section{既存文献との関係}

\textbf{$p$-adic物理学。} Volovich~\cite{Volovich1987}の$p$-adic時空、Dragovich~\cite{Dragovich2009}の$p$-adic宇宙論が数学的基盤を提供。我々は$\Spec(\mathbb{Z})$を通じたアデール的設定に拡張。

\textbf{NCG標準模型。} Chamseddine-Connes~\cite{ChamseddineConnes1997}がスペクトル作用原理を導入、CCM~\cite{CCM2007}がSM全体を導出。$\alpha$公式はスペクトル作用のモーメントが$\zeta$値で決まれば導出可能。

\textbf{ホログラフィックダークエネルギー。} Li~\cite{Li2004}の$\rho_\Lambda \propto M_P^2 L^{-2}$と同形だが、係数$\zeta(-3) = 1/120$を提供。

\textbf{物理学における$\zeta$関数。} Hawking~\cite{Hawking1977}, Elizalde~\cite{Elizalde2012}, Schumayer-Hutchinson~\cite{Schumayer2011}, Manin~\cite{Manin1989}の仕事を\emph{予測的}枠組みに拡張。

\textbf{ミューオン$g-2$。} Fermilab測定~\cite{Muong2_2021}は通常BSM物理（SUSY等~\cite{Athron2021}）の証拠と解釈される。我々は算術的代替を提案：異常は$\zeta(-5)$補正であり、新粒子シグナルではない。

\textbf{ニュートリノ質量。} $\Delta m^2_{32}/\Delta m^2_{21} \approx 33 = \pi(137)$の接続は文献に存在しない新規観察。

\section{反証条件}

仮説は反証可能：
\begin{itemize}
\item CMB-S4が$r$を$5\sigma$精度で測定し$\pi^2/720$と非両立：予測5反証。
\item DESI/Euclidが$w(z) = -1.000 \pm 0.005$を離散特徴なしで制約：予測7反証。
\item Planck再解析が$\ell \approx 328, 390, 474$に余剰なし：予測6反証。
\item 格子QCD+分散関係が$g-2$の不一致を解消し、最終的$\Delta a_\mu$が$251 \times 10^{-11}$から有意に逸脱：予測4弱体化。
\end{itemize}

\section{結論}

仮説「時空$= \Spec(\mathbb{Z})$」を10の定量的テストに付した。5つが既存データと比較可能であり、全て整合（$0.00002\%$から1桁まで）。5つが近未来テスト可能。

最も印象的な特徴は：(1)~自由パラメータゼロ。(2)~同じ算術量が無関係な物理領域に出現する領域横断的整合性。(3)~標準的アプローチからの改善（ダークエネルギーで$10^{121}$桁、$\alpha$は標準模型に対応物なし）。(4)~ニュートリノ質量比$\approx \pi(1/\alpha)$、ミューオン異常$\approx (1/|\zeta(-5)|-1) \times 10^{-11}$という既存の枠組みが予期しない新規接続。

特にミューオン$g-2$（予測4）は、整数値$251$がチューニング不能な形で実験中心値と完全一致するという点で、最も「こじつけ」から遠い証拠である。

最も即座のテスト——Planck CMB再解析（予測6）とDESI/Euclid $w(z)$測定（予測7）——は新しい実験を必要とせず、既存または近日公開のデータの新しい解析のみを要する。

\section*{謝辞}
本研究はWright Brothers共同研究グループによる独立研究として遂行された。

\begin{thebibliography}{30}
\bibitem{Connes1994} A.~Connes, \emph{Noncommutative Geometry}, Academic Press (1994).
\bibitem{Connes1996} A.~Connes, Commun. Math. Phys. \textbf{182}, 155 (1996).
\bibitem{ConnesMarcolli2008} A.~Connes and M.~Marcolli, \emph{NCG, Quantum Fields and Motives}, AMS (2008).
\bibitem{CCM2007} A.~H.~Chamseddine et al., Adv. Theor. Math. Phys. \textbf{11}, 991 (2007).
\bibitem{ChamseddineConnes1997} A.~H.~Chamseddine and A.~Connes, Commun. Math. Phys. \textbf{186}, 731 (1997).
\bibitem{Volovich1987} I.~V.~Volovich, Class. Quantum Grav. \textbf{4}, L83 (1987).
\bibitem{Dragovich2009} B.~Dragovich et al., $p$-Adic Numbers Ultrametric Anal. Appl. \textbf{1}, 1 (2009).
\bibitem{Brekke1993} L.~Brekke and P.~G.~O.~Freund, Phys. Rep. \textbf{233}, 1 (1993).
\bibitem{Li2004} M.~Li, Phys. Lett. B \textbf{603}, 1 (2004).
\bibitem{Hawking1977} S.~W.~Hawking, Commun. Math. Phys. \textbf{55}, 133 (1977).
\bibitem{Elizalde2012} E.~Elizalde, \emph{Spectral Zeta Functions}, Springer (2012).
\bibitem{Schumayer2011} D.~Schumayer et al., Rev. Mod. Phys. \textbf{83}, 307 (2011).
\bibitem{Manin1989} Yu.~I.~Manin, in \emph{Conformal Invariance and String Theory}, Academic Press (1989).
\bibitem{Weinberg1989} S.~Weinberg, Rev. Mod. Phys. \textbf{61}, 1 (1989).
\bibitem{Mohr2016} P.~J.~Mohr et al., Rev. Mod. Phys. \textbf{88}, 035009 (2016).
\bibitem{Esteban2020} I.~Esteban et al., JHEP \textbf{2020}, 178 (2020).
\bibitem{Muong2_2021} Muon $g$-2 Collab., Phys. Rev. Lett. \textbf{126}, 141801 (2021).
\bibitem{Athron2021} P.~Athron et al., JHEP \textbf{2021}, 80 (2021).
\bibitem{BICEP2021} BICEP/Keck Collab., Phys. Rev. Lett. \textbf{127}, 151301 (2021).
\bibitem{CMBS4} CMB-S4 Collab., arXiv:1610.02743 (2016).
\bibitem{Planck2020} Planck Collab., A\&A \textbf{641}, A6 (2020).
\bibitem{DESI2024} DESI Collab., arXiv:2404.03002 (2024).
\bibitem{Euclid2024} Euclid Collab., A\&A \textbf{681}, A67 (2024).
\bibitem{Sen2012} A.~Sen, Gen. Rel. Grav. \textbf{44}, 1207 (2012).
\bibitem{FermiLAT2010} Fermi-LAT Collab., Science \textbf{328}, 725 (2010).
\bibitem{BostConnes1995} J.-B.~Bost and A.~Connes, Selecta Math. \textbf{1}, 411 (1995).
\end{thebibliography}

\end{document}
